% =============================================================================
% Presentación Beamer P05 – Análisis de Imagen Digital y Fractografía SEM
% Técnicas de Medida – Ingeniería Aeronáutica
% Compilar con: latexmk -xelatex -interaction=nonstopmode presentacion_P05.tex
% =============================================================================
\documentclass[aspectratio=169,12pt]{beamer}

\usetheme{Madrid}
\usecolortheme{seahorse}
\usefonttheme{professionalfonts}

% ---- Encoding & Language ----
\usepackage{fontspec}
\usepackage{csquotes}
\usepackage[spanish]{babel}
\addto\captionsspanish{\renewcommand{\tablename}{Tabla}}

% ---- Math & Physics ----
\usepackage{amsmath,amssymb}
\usepackage{mathtools}
\usepackage{siunitx}
\sisetup{output-decimal-marker={,}}

% ---- Graphics ----
\usepackage{graphicx}
\graphicspath{{../../../Practicas/P05_Analisis_Imagen/data/images/}
              {../../../files/p5/}}

% ---- TikZ ----
\usepackage{tikz}
\usetikzlibrary{babel}
\usetikzlibrary{arrows.meta,positioning,calc,shapes.geometric,
                decorations.markings,decorations.pathmorphing,patterns}

% ---- Tables ----
\usepackage{booktabs}
\usepackage{colortbl}
\usepackage{tabularx}
\usepackage{array}

% ---- Code ----
\usepackage{listings}
\lstset{
  basicstyle=\ttfamily\tiny,
  keywordstyle=\color{blue!80}\bfseries,
  commentstyle=\color{green!50!black},
  backgroundcolor=\color{gray!10},
  frame=single,
  breaklines=true
}

% ---- Custom Colors ----
\definecolor{aeroblue}{RGB}{0,51,102}
\definecolor{aerored}{RGB}{153,0,0}
\definecolor{aerogreen}{RGB}{0,102,51}
\definecolor{aeroyellow}{RGB}{204,153,0}
\definecolor{aerogreendark}{RGB}{0,80,40}
\definecolor{aerogreenlight}{RGB}{0,140,70}

% ---- GREEN THEME ----
\setbeamercolor{frametitle}{bg=aerogreen,fg=white}
\setbeamercolor{title}{bg=aerogreen,fg=white}
\setbeamercolor{structure}{fg=aerogreen}
\setbeamercolor{palette primary}{bg=aerogreen,fg=white}
\setbeamercolor{palette secondary}{bg=aerogreenlight,fg=white}
\setbeamercolor{palette tertiary}{bg=aerogreendark,fg=white}
\setbeamercolor{palette quaternary}{bg=aerogreendark,fg=white}
\setbeamercolor{section in toc}{fg=aerogreen}
\setbeamercolor{subsection in toc}{fg=aerogreendark}
\setbeamercolor{block title}{bg=aerogreen,fg=white}
\setbeamercolor{block body}{bg=aerogreen!10,fg=black}
\setbeamercolor{block title alerted}{bg=aerored,fg=white}
\setbeamercolor{block body alerted}{bg=aerored!10,fg=black}
\setbeamercolor{block title example}{bg=aerogreenlight,fg=white}
\setbeamercolor{block body example}{bg=aerogreenlight!10,fg=black}
\setbeamercolor{item}{fg=aerogreen}
\setbeamercolor{itemize item}{fg=aerogreen}
\setbeamercolor{enumerate item}{fg=aerogreen}

% ---- Title Info ----
\title[P05 -- Análisis de Imagen]{P05: Análisis de Imagen Digital\\y Fractografía SEM}
\subtitle{Técnicas de Medida -- Ingeniería Aeronáutica}
\author{Laboratorio de Técnicas de Medida}
\date{}

% =============================================================================
\begin{document}

\begin{frame}
	\titlepage
\end{frame}

\begin{frame}{Índice}
	\tableofcontents
\end{frame}

% =============================================================================
\section{Introducción}
% =============================================================================

\begin{frame}{¿Qué es la Fractografía SEM?}
	\begin{columns}
		\begin{column}{0.55\textwidth}
			\begin{itemize}\small
				\item Estudio de superficies de fractura con
				      \textbf{Microscopía Electrónica de Barrido}
				\item Resolución $\sim\!\SI{1}{\nano\m}$,
				      gran profundidad de campo
				\item Identifica el \textbf{modo de falla}:
				      dúctil (hoyuelos), frágil (clivaje),
				      fatiga (estrías)
				\item Complemento: \textbf{procesamiento digital}
				      para cuantificación automatizada
			\end{itemize}
			\vspace{0.1em}
			\begin{alertblock}{Objetivo}
				\scriptsize Pipeline de 6 etapas sobre micrografías SEM
				+ mediciones dimensionales
			\end{alertblock}
		\end{column}
		\begin{column}{0.42\textwidth}
			\begin{tikzpicture}[scale=0.55,>=Stealth]
				% SEM column
				\draw[fill=gray!20,draw=aeroblue,thick]
				(-0.6,3.2) rectangle (0.6,4.2);
				\node[aeroblue,font=\tiny\bfseries] at (0,3.7) {Cañón $e^-$};
				% Lens
				\draw[aeroblue,thick] (-0.8,2.6) arc (180:360:0.8 and 0.12);
				\draw[aeroblue,thick] (-0.8,2.6) arc (180:0:0.8 and 0.12);
				% Beam
				\draw[blue!50,line width=1.5pt] (0,3.2) -- (0,0.3);
				% Sample
				\draw[fill=orange!15,draw=black,thick]
				(-1.2,0) -- (1.2,0) -- (0.9,-0.4) -- (-0.9,-0.4) -- cycle;
				\node[font=\tiny] at (0,-0.7) {Fractura};
				% Detector
				\draw[fill=aerogreen!20,draw=aerogreen,thick]
				(1.6,1.0) circle (0.3);
				\node[aerogreen,font=\tiny] at (1.6,0.5) {SE};
				\draw[->,dashed,gray] (0.2,0.2) -- (1.3,0.8);
				% Arrow to processing
				\draw[aeroblue,ultra thick,->] (2.3,1.0) -- (3.5,1.0);
				\node[above,font=\tiny,aeroblue] at (2.9,1.1) {Imagen};
				% Pipeline
				\draw[fill=aeroblue!10,draw=aeroblue,thick,rounded corners=2pt]
				(3.5,0.3) rectangle (5.0,1.7);
				\node[aeroblue,font=\tiny\bfseries,align=center] at (4.25,1.0)
				{CLAHE\\Sobel\\Canny\\Otsu};
			\end{tikzpicture}
		\end{column}
	\end{columns}
\end{frame}

\begin{frame}{Aplicaciones en Ingeniería Aeronáutica}
	\begin{columns}
		\begin{column}{0.5\textwidth}
			\begin{enumerate}\small
				\item Diagnóstico de causa raíz de fallas en servicio
				\item Control de calidad de aleaciones aeronáuticas
				\item Verificación de tratamientos superficiales
				\item Estimación de vida a fatiga ($da/dN$ por estrías)
				\item Validación de modelos de mecánica de fractura
			\end{enumerate}
			\vspace{0.3em}
			\begin{exampleblock}{Dato clave}
				La \textbf{fatiga} causa $>60\%$ de las fallas
				estructurales en aeronáutica
			\end{exampleblock}
		\end{column}
		\begin{column}{0.47\textwidth}
			\begin{alertblock}{¿Por qué automatizar?}
				\begin{itemize}\small
					\item Medición manual es lenta y subjetiva
					\item Pipeline automatizado: reproducible,
					      trazable, cuantitativo
					\item Estadísticas sobre miles de regiones
					\item Comparación objetiva entre muestras
				\end{itemize}
			\end{alertblock}
		\end{column}
	\end{columns}
\end{frame}

\begin{frame}{Objetivos de Aprendizaje}
	\begin{enumerate}
		\item Comprender el principio del SEM y los parámetros de
		      adquisición (EHT, WD, detector, escala)
		\item Aplicar la ecualización adaptativa CLAHE para mejorar
		      el contraste local
		\item Calcular bordes con los operadores Sobel y Canny
		\item Segmentar con la umbralización óptima de Otsu
		\item Refinar máscaras con operaciones morfológicas (opening)
		\item Etiquetar regiones y extraer descriptores de forma
		\item Calibrar la escala espacial (\si{px}/\si{\micro\m})
		      y medir dimensiones
		\item Diagnosticar modos de fractura a partir de evidencia
		      cuantitativa
	\end{enumerate}
\end{frame}

% =============================================================================
\section{Fundamentos Teóricos}
% =============================================================================

\begin{frame}{Microscopía Electrónica de Barrido (SEM)}
	\begin{columns}
		\begin{column}{0.5\textwidth}
			\begin{tikzpicture}[scale=0.6,>=Stealth]
				% Column
				\draw[fill=gray!15,draw=black,thick] (-0.7,4.5) rectangle (0.7,5.5);
				\node[font=\scriptsize] at (0,5.0) {Cañón $e^-$};
				% Condenser
				\draw[thick,aeroblue] (-1.0,3.8) arc (180:360:1.0 and 0.15);
				\draw[thick,aeroblue] (-1.0,3.8) arc (180:0:1.0 and 0.15);
				\node[right,font=\tiny,aeroblue] at (1.1,3.8) {Condensadora};
				% Scan coils
				\draw[fill=aeroyellow!30,draw=black]
				(-0.4,3.0) rectangle (-0.15,3.4);
				\draw[fill=aeroyellow!30,draw=black]
				(0.15,3.0) rectangle (0.4,3.4);
				\node[font=\tiny] at (0,2.7) {Bobinas};
				% Objective
				\draw[thick,aeroblue] (-0.8,2.2) arc (180:360:0.8 and 0.12);
				\draw[thick,aeroblue] (-0.8,2.2) arc (180:0:0.8 and 0.12);
				\node[right,font=\tiny,aeroblue] at (0.9,2.2) {Objetivo};
				% Beam
				\draw[blue!60,line width=1.5pt] (0,4.5) -- (0,0.3);
				% Sample
				\draw[fill=orange!15,draw=black,thick]
				(-1.5,0) -- (1.5,0) -- (1.2,-0.5) -- (-1.2,-0.5) -- cycle;
				\node[font=\scriptsize] at (0,-0.8) {Muestra};
				% SE detector
				\draw[fill=aerogreen!20,draw=aerogreen,thick]
				(2.2,1.2) ellipse (0.35 and 0.25);
				\node[aerogreen,font=\tiny] at (2.2,0.7) {Detector SE};
				\draw[->,dashed,aerogreen] (0.3,0.2) -- (1.9,1.0);
			\end{tikzpicture}
		\end{column}
		\begin{column}{0.47\textwidth}
			\begin{block}{Señales del SEM}
				\begin{itemize}\small
					\item \textbf{SE} (secundarios):
					      topografía superficial
					\item \textbf{BSE} (retrodispersados):
					      contraste composicional
					\item \textbf{Rayos~X}: análisis elemental
					      (EDS/EDX)
				\end{itemize}
			\end{block}
			\vspace{0.3em}
			\begin{exampleblock}{Ventajas vs.\ óptico}
				\begin{itemize}\scriptsize
					\item Resolución: $\sim\!\SI{1}{\nano\m}$
					      (vs.\ $\sim\!\SI{200}{\nano\m}$)
					\item Profundidad de campo $\gg$
					\item Imagen 3D de superficies rugosas
				\end{itemize}
			\end{exampleblock}
		\end{column}
	\end{columns}
\end{frame}

\begin{frame}{Parámetros de Adquisición SEM}
	\centering
	\small
	\begin{tabular}{lll}
		\toprule
		\textbf{Parámetro}       & \textbf{Unidad}
		                         & \textbf{Efecto}                                 \\
		\midrule
		EHT (voltaje)            & \si{\kilo\volt}
		                         & $\uparrow$ $\to$ mayor penetración y resolución \\
		WD (dist.\ trabajo)      & \si{\mm}
		                         & $\downarrow$ $\to$ mejor resolución             \\
		Mag (magnificación)      & $\times$
		                         & Define campo de visión                          \\
		Signal A (detector)      & SE / BSE
		                         & SE: topografía; BSE: composición                \\
		\rowcolor{aeroyellow!20}
		\textbf{Barra de escala} & \si{\micro\m}
		                         & \textbf{Fundamental para mediciones}            \\
		\bottomrule
	\end{tabular}

	\vspace{0.5cm}
	\begin{tikzpicture}
		\node[draw=aerogreen,thick,rounded corners,fill=aerogreen!10,
			text width=10cm,align=center,font=\small] {
			La barra de escala define la relación
			\textbf{píxeles $\leftrightarrow$ \si{\micro\m}}.\\
			Sin ella, no hay medición cuantitativa posible.
		};
	\end{tikzpicture}
\end{frame}

\begin{frame}{Mecanismos de Fractura}
	\begin{columns}
		\begin{column}{0.32\textwidth}
			\begin{block}{\centering Dúctil}
				\centering
				\begin{tikzpicture}[scale=0.55]
					\draw[thick] (0,0) rectangle (2.8,2.2);
					\foreach \x/\y/\r in {0.4/1.0/0.2, 1.0/0.6/0.25,
							1.4/1.4/0.18, 2.0/0.9/0.3, 2.4/1.6/0.15,
							0.7/1.7/0.12, 1.8/0.3/0.2} {
							\draw[gray,fill=gray!20] (\x,\y)
							ellipse ({\r} and {\r*0.7});
						}
				\end{tikzpicture}\\
				\scriptsize \textbf{Hoyuelos} (dimples)\\
				$G_\mathrm{dúctil} \gg G_\mathrm{frágil}$
			\end{block}
		\end{column}
		\begin{column}{0.32\textwidth}
			\begin{block}{\centering Frágil}
				\centering
				\begin{tikzpicture}[scale=0.55]
					\draw[thick] (0,0) rectangle (2.8,2.2);
					\draw[gray] (0.2,0.1) -- (1.0,1.6) -- (1.4,0.4) -- cycle;
					\draw[gray] (1.0,1.6) -- (2.0,1.8) -- (1.4,0.4) -- cycle;
					\draw[gray] (1.4,0.4) -- (2.0,1.8) -- (2.6,0.6) -- cycle;
					\draw[aerored,thin,->] (0.7,0.8) -- (1.1,1.0);
					\draw[aerored,thin,->] (1.5,1.0) -- (1.9,1.3);
				\end{tikzpicture}\\
				\scriptsize \textbf{Clivaje} + marcas de río\\
				Planos cristalográficos
			\end{block}
		\end{column}
		\begin{column}{0.32\textwidth}
			\begin{block}{\centering Fatiga}
				\centering
				\begin{tikzpicture}[scale=0.55]
					\draw[thick] (0,0) rectangle (2.8,2.2);
					\foreach \y in {0.25,0.5,0.75,1.0,1.25,1.5,1.75,2.0} {
							\draw[gray,thin] (0.15,\y)
							.. controls (0.8,{\y+0.1}) and (2.0,{\y-0.08})
							.. (2.65,\y);
						}
					\fill[aerored] (0.15,1.1) circle (0.06);
					\draw[aerored,->] (0.15,1.1) -- (0.6,1.1);
				\end{tikzpicture}\\
				\scriptsize \textbf{Estrías} de fatiga\\
				$da/dN = C(\Delta K)^m$
			\end{block}
		\end{column}
	\end{columns}

	\vspace{0.4cm}
	\centering
	\begin{tikzpicture}
		\node[draw=aerored,thick,rounded corners,fill=aerored!8,
			text width=10cm,align=center,font=\small] {
			Nuestras 5 micrografías muestran \textbf{superficies de
				fractura por fatiga} con estrías, marcas de playa y
			zonas de fractura final.
		};
	\end{tikzpicture}
\end{frame}

\begin{frame}{Ley de Paris -- Propagación de Grietas por Fatiga}
	\begin{columns}
		\begin{column}{0.5\textwidth}
			\textbf{Tasa de crecimiento:}
			\begin{equation*}
				\boxed{\frac{da}{dN} = C\,(\Delta K)^m}
			\end{equation*}

			\begin{itemize}\small
				\item $a$: longitud de grieta
				\item $N$: número de ciclos
				\item $\Delta K$: rango del factor de intensidad
				\item $C$, $m$: constantes del material
			\end{itemize}

			\vspace{0.3em}
			\begin{alertblock}{Conexión con estrías}
				Cada estría $\approx$ 1 ciclo de carga\\
				Espaciado entre estrías $\approx da/dN$
			\end{alertblock}
		\end{column}
		\begin{column}{0.47\textwidth}
			\begin{tikzpicture}[scale=0.65]
				\draw[->] (0,0) -- (5.5,0)
				node[right,font=\tiny] {$\log(\Delta K)$};
				\draw[->] (0,0) -- (0,4.5)
				node[above,font=\tiny] {$\log(da/dN)$};
				% Three regions
				\fill[aeroblue!10] (0,0) rectangle (1.5,4.5);
				\fill[aerogreen!10] (1.5,0) rectangle (3.8,4.5);
				\fill[aerored!8] (3.8,0) rectangle (5.5,4.5);
				\node[aeroblue,font=\tiny,rotate=90] at (0.75,2.2)
				{Región I};
				\node[aerogreen,font=\tiny,rotate=90] at (2.65,2.2)
				{Región II (Paris)};
				\node[aerored,font=\tiny,rotate=90] at (4.65,2.2)
				{Región III};
				% Curve
				\draw[aerored,very thick]
				(0.8,0.3) .. controls (1.2,0.6) .. (1.5,1.0)
				-- (3.8,3.2)
				.. controls (4.2,3.8) .. (4.8,4.3);
				% Slope label
				\node[aerogreen,font=\small\bfseries] at (3.0,1.8)
				{pend.\ $= m$};
				% Threshold
				\draw[aerogreen,dashed] (1.5,0) -- (1.5,4.5);
				\node[below,font=\tiny] at (1.5,0) {$\Delta K_\mathrm{th}$};
			\end{tikzpicture}
		\end{column}
	\end{columns}
\end{frame}

% =============================================================================
\section{Pipeline de Procesamiento}
% =============================================================================

\begin{frame}{Pipeline de 6 Etapas}
	\centering
	\begin{tikzpicture}[node distance=0.6cm,>=Stealth,
			block/.style={rectangle,draw=aerogreen,thick,fill=aerogreen!8,
					text width=4.0cm,align=center,minimum height=0.6cm,
					font=\scriptsize,rounded corners=2pt}]
		\node[block] (load) {\textbf{1.} Carga en gris $[0,1]$};
		\node[block,below=of load] (clahe)
		{\textbf{2.} CLAHE (\texttt{clip}=0,03)};
		\node[block,below=of clahe] (sobel)
		{\textbf{3a.} Sobel (gradiente)};
		\node[block,right=2cm of sobel] (canny)
		{\textbf{3b.} Canny ($\sigma=2$)};
		\node[block,below=of sobel] (otsu)
		{\textbf{4.} Umbralización Otsu};
		\node[block,below=of otsu] (morph)
		{\textbf{5.} Opening (disk=2)};
		\node[block,below=of morph,draw=aerored,fill=aerored!8]
		(label) {\textbf{6.} Etiquetado $\to$ $N$ regiones};
		% Arrows
		\foreach \a/\b in {load/clahe,clahe/sobel,sobel/otsu,
				otsu/morph,morph/label} {
				\draw[->,aerogreen,thick] (\a) -- (\b);
			}
		\draw[->,aerogreen,thick] (clahe.east) -| (canny.north);
		% Side annotations
		\node[right=0.5cm of load,font=\tiny,text width=3cm]
		{float, normalizado};
		\node[right=0.5cm of clahe,font=\tiny,text width=3cm]
		{Contraste local mejorado};
		\node[right=0.5cm of otsu,font=\tiny,text width=3cm]
		{Máscara binaria automática};
		\node[right=0.5cm of morph,font=\tiny,text width=3cm]
		{Ruido eliminado};
		\node[right=0.5cm of label,font=\tiny,text width=3cm,aerored]
		{Descriptores de forma};
	\end{tikzpicture}
\end{frame}

\begin{frame}{Etapa 2: CLAHE -- Ecualización Adaptativa}
	\begin{columns}
		\begin{column}{0.5\textwidth}
			\textbf{Problema:} iluminación desigual en SEM

			\vspace{0.3em}
			\textbf{Solución CLAHE:}
			\begin{equation*}
				I_\mathrm{CLAHE}(x,y)
				= T_\mathrm{tile}[I(x,y)]
			\end{equation*}

			\begin{itemize}\small
				\item Divide imagen en \emph{tiles}
				\item Ecualiza cada tile localmente
				\item \texttt{clip\_limit}: limita la amplificación
				      del contraste (evita ruido)
				\item Interpolación bilineal entre tiles
			\end{itemize}

			\vspace{0.3em}
			\begin{block}{Parámetro}
				\texttt{clip\_limit = 0,03} (normalizado)
			\end{block}
		\end{column}
		\begin{column}{0.47\textwidth}
			\includegraphics[width=\textwidth]{fatigue3_01_clahe.png}
			\centering
			\scriptsize Resultado CLAHE sobre \texttt{fatigue3.png}
		\end{column}
	\end{columns}
\end{frame}

\begin{frame}{Etapa 3a: Operador de Sobel}
	\begin{columns}
		\begin{column}{0.5\textwidth}
			\textbf{Kernels de gradiente:}
			\begin{equation*}
				G_x = \begin{bmatrix}
					-1 & 0 & +1 \\ -2 & 0 & +2 \\ -1 & 0 & +1
				\end{bmatrix} \ast I
			\end{equation*}
			\begin{equation*}
				G_y = \begin{bmatrix}
					+1 & +2 & +1 \\ 0 & 0 & 0 \\ -1 & -2 & -1
				\end{bmatrix} \ast I
			\end{equation*}

			\textbf{Magnitud del gradiente:}
			\begin{equation*}
				\boxed{|\nabla I| = \sqrt{G_x^2 + G_y^2}}
			\end{equation*}

			$\to$ Valores altos = bordes
		\end{column}
		\begin{column}{0.47\textwidth}
			\includegraphics[width=\textwidth]{fatigue3_02_sobel.png}
			\centering
			\scriptsize Filtro Sobel sobre \texttt{fatigue3.png}
		\end{column}
	\end{columns}
\end{frame}

\begin{frame}{Etapa 3b: Detector de Bordes de Canny}
	\begin{columns}
		\begin{column}{0.5\textwidth}
			\textbf{4 pasos del algoritmo:}
			\begin{enumerate}\small
				\item \textbf{Suavizado Gaussiano} ($\sigma = 2{,}0$)
				      \begin{equation*}
					      G_\sigma = \frac{1}{2\pi\sigma^2}
					      e^{-\frac{x^2+y^2}{2\sigma^2}}
				      \end{equation*}
				\item \textbf{Gradiente} $|\nabla I|$ y orientación
				      $\theta$
				\item \textbf{Supresión no máxima}:\\
				      solo máximos locales en dirección de $\theta$
				\item \textbf{Histéresis}:\\
				      $T_\mathrm{high}$ confirma, $T_\mathrm{low}$
				      extiende
			\end{enumerate}
		\end{column}
		\begin{column}{0.47\textwidth}
			\includegraphics[width=0.92\textwidth]{fatigue3_03_canny.png}
			\centering
			\scriptsize Bordes Canny ($\sigma=2$) sobre
			\texttt{fatigue3.png}

			\vspace{0.2em}
			\begin{exampleblock}{Ventaja}
				\scriptsize Bordes de 1~px, óptimos para
				medir contornos de estrías
			\end{exampleblock}
		\end{column}
	\end{columns}
\end{frame}

\begin{frame}{Etapa 4: Umbralización de Otsu}
	\begin{columns}
		\begin{column}{0.5\textwidth}
			\textbf{Minimiza varianza intra-clase}\\
			(equivalente a maximizar inter-clase):

			\begin{equation*}
				\sigma_B^2(t) = \omega_0\,\omega_1\,
				(\mu_0 - \mu_1)^2
			\end{equation*}

			\begin{equation*}
				\boxed{t^* = \arg\max_{t}\;\sigma_B^2(t)}
			\end{equation*}

			\begin{itemize}\small
				\item $\omega_0$, $\omega_1$: probabilidades de clase
				\item $\mu_0$, $\mu_1$: medias de clase
				\item Resultado: máscara binaria automática
				\item No requiere parámetros del usuario
			\end{itemize}
		\end{column}
		\begin{column}{0.47\textwidth}
			\includegraphics[width=0.92\textwidth]{fatigue3_04_otsu.png}
			\centering
			\scriptsize Umbralización Otsu sobre
			\texttt{fatigue3.png}

			\vspace{0.2em}
			\begin{alertblock}{Limitación}
				\scriptsize Requiere histograma bimodal.
				Se aplica tras CLAHE.
			\end{alertblock}
		\end{column}
	\end{columns}
\end{frame}

\begin{frame}{Etapa 5: Morfología -- Opening}
	\begin{columns}
		\begin{column}{0.5\textwidth}
			\textbf{Erosión + Dilatación:}
			\begin{equation*}
				\boxed{I \circ SE = (I \ominus SE) \oplus SE}
			\end{equation*}

			\vspace{-0.3em}
			\begin{itemize}\small
				\item \textbf{Erosión:}
				      $(I \ominus SE) = \min_{SE}\,I$
				\item \textbf{Dilatación:}
				      $(I \oplus SE) = \max_{SE}\,I$
				\item Opening elimina objetos $<$ SE
				      (disco radio~2, $\pi(2)^2 \approx 12{,}6$~px$^2$)
			\end{itemize}

			\vspace{0.2em}
			\begin{block}{Propiedad}
				$I \circ SE \subseteq I$ (nunca agrega píxeles)
			\end{block}
		\end{column}
		\begin{column}{0.47\textwidth}
			\includegraphics[width=0.92\textwidth]{fatigue3_05_opening.png}
			\centering
			\scriptsize Opening (disk=2) sobre
			\texttt{fatigue3.png}

			\vspace{0.2em}
			\begin{exampleblock}{Resultado}
				\scriptsize Ruido eliminado, regiones
				preservadas. Listo para etiquetado.
			\end{exampleblock}
		\end{column}
	\end{columns}
\end{frame}

\begin{frame}{Etapa 6: Etiquetado -- Descriptores de Región}
	\begin{columns}
		\begin{column}{0.5\textwidth}
			\textbf{Componentes conexos:} cada grupo de píxeles
			``1'' conectados $\to$ etiqueta única
			$\ell \in \{1, \ldots, N\}$

			\vspace{0.5em}
			\textbf{Descriptores por región:}
			\centering\small
			\begin{tabular}{ll}
				\toprule
				\textbf{Descriptor} & \textbf{Significado}    \\
				\midrule
				Área ($A_r$)        & Tamaño (px$^2$)         \\
				Perímetro ($P_r$)   & Longitud borde (px)     \\
				Excentricidad ($e$) & Alargamiento ($0$--$1$) \\
				Eje mayor           & Dimensión máxima (px)   \\
				\bottomrule
			\end{tabular}
		\end{column}
		\begin{column}{0.47\textwidth}
			\begin{block}{Uso}
				Clasificar regiones como estrías (alta $e$)
				vs.\ hoyuelos (baja $e$)
			\end{block}

			\vspace{0.3em}
			\begin{tikzpicture}[scale=0.55]
				\draw[thick] (0,0) rectangle (4,3);
				% Labelled regions
				\fill[aerogreen!40] (0.3,0.5) ellipse (0.5 and 0.3);
				\node[font=\tiny\bfseries,white] at (0.3,0.5) {1};
				\fill[aeroblue!40] (1.5,1.8) ellipse (0.8 and 0.2);
				\node[font=\tiny\bfseries,white] at (1.5,1.8) {2};
				\fill[aerored!40] (3.0,2.2) ellipse (0.4 and 0.5);
				\node[font=\tiny\bfseries,white] at (3.0,2.2) {3};
				\fill[aeroyellow!60] (2.5,0.6) ellipse (0.6 and 0.25);
				\node[font=\tiny\bfseries] at (2.5,0.6) {4};
				\node[below,font=\scriptsize] at (2,-0.2)
				{Ejemplo: 4 regiones etiquetadas};
			\end{tikzpicture}
		\end{column}
	\end{columns}
\end{frame}

\begin{frame}{Etapa 6: Resultados del Etiquetado}
	\centering
	\textbf{Regiones detectadas por imagen:}

	\vspace{0.5em}
	\begin{tabular}{l r l}
		\toprule
		\textbf{Imagen} & $\boldsymbol{N}$ & \textbf{Observación} \\
		\midrule
		fatigue3        & 418              & Densidad media       \\
		fatigue4        & 153              & Buen contraste       \\
		\rowcolor{aeroyellow!20}
		fatigue5        & 2066             & Alta resolución      \\
		\rowcolor{aerored!15}
		fatigue6        & 2897             & Máxima complejidad   \\
		fatigue7        & 549              & Textura fina         \\
		\bottomrule
	\end{tabular}

	\vspace{0.6em}
	\begin{columns}
		\begin{column}{0.48\textwidth}
			\begin{alertblock}{Observación}
				$N$ depende fuertemente de la
				resolución/magnificación de la imagen.
			\end{alertblock}
		\end{column}
		\begin{column}{0.48\textwidth}
			\begin{exampleblock}{Consistencia}
				El pipeline es determinista: mismos
				parámetros $\to$ mismos $N$.
			\end{exampleblock}
		\end{column}
	\end{columns}
\end{frame}

% =============================================================================
\section{Resultados: \texttt{fatigue3.png}}
% =============================================================================

\begin{frame}{Resultado Completo: \texttt{fatigue3.png}}
	\centering
	\includegraphics[width=\textwidth,height=0.78\textheight,keepaspectratio]{fatigue3_00_resumen.png}

	\vspace{0.1cm}
	\scriptsize Figura compuesta: Histograma $\to$ CLAHE $\to$ Sobel $\to$
	Canny $\to$ Otsu $\to$ Opening. \textbf{418 regiones} detectadas.
\end{frame}

\begin{frame}{Interpretación del Pipeline -- \texttt{fatigue3.png}}
	\begin{columns}
		\begin{column}{0.48\textwidth}
			\begin{enumerate}\scriptsize
				\item \textbf{Histograma:} concentrado en tonos
				      medio-bajos $\to$ bajo contraste
				\item \textbf{CLAHE:} mejora detalles finos
				      (estrías, bordes) sin saturar
				\item \textbf{Sobel:} revela textura;
				      alto gradiente = transiciones
				\item \textbf{Canny:} contornos nítidos (1~px)
				\item \textbf{Otsu:} segmentación automática
				      (claras vs.\ oscuras)
				\item \textbf{Opening:} limpieza $\to$
				      \textbf{418 regiones}
			\end{enumerate}
		\end{column}
		\begin{column}{0.5\textwidth}
			\begin{block}{Diagnóstico preliminar}
				\scriptsize
				\begin{itemize}
					\item Estrías de fatiga visibles (líneas
					      paralelas)
					\item Dirección de propagación identificable
					\item Zona de fractura final (rasgos irregulares)
				\end{itemize}
			\end{block}

			\vspace{0.2em}
			\begin{exampleblock}{Mediciones necesarias}
				\scriptsize
				\begin{itemize}
					\item Espaciado de estrías $\to$ $da/dN$
					\item Tamaño de regiones $\to$ distribución
					\item Excentricidad $\to$ clasificación
				\end{itemize}
			\end{exampleblock}
		\end{column}
	\end{columns}
\end{frame}

% =============================================================================
\section{Calibración y Medición}
% =============================================================================

\begin{frame}{Calibración de Escala}
	\begin{columns}
		\begin{column}{0.5\textwidth}
			\textbf{Procedimiento:}
			\begin{enumerate}\small
				\item Clic en extremos $P_1$, $P_2$ de la barra
				\item Calcular distancia en píxeles:
				      \begin{equation*}
					      d_\mathrm{px}
					      = \sqrt{(x_2-x_1)^2 + (y_2-y_1)^2}
				      \end{equation*}
				\item Ingresar distancia real
				      $d_\mathrm{real}$ (\si{\micro\m})
				\item Calcular escala:
				      \begin{equation*}
					      \boxed{s = \frac{d_\mathrm{px}}{d_\mathrm{real}}
						      \;\left[\frac{\mathrm{px}}{\si{\micro\m}}\right]}
				      \end{equation*}
			\end{enumerate}
		\end{column}
		\begin{column}{0.47\textwidth}
			\textbf{Error de calibración:}
			\begin{equation*}
				\varepsilon_s = \frac{2n}{d_\mathrm{px}} \times 100\%
			\end{equation*}

			\begin{exampleblock}{Ejemplo}
				\scriptsize
				Si la barra mide 200~px y cometemos
				$\pm$5~px de error:
				\[
					\varepsilon_s = \frac{2 \times 5}{200}
					\times 100\% = 5\%
				\]
				$\to$ todas las mediciones heredan este 5\%
			\end{exampleblock}

			\vspace{0.3em}
			\begin{tikzpicture}[scale=0.6]
				\draw[thick] (0,0) -- (4,0);
				\fill[aerored] (0,0) circle (0.08);
				\fill[aerored] (4,0) circle (0.08);
				\node[above,font=\tiny] at (0,0.1) {$P_1$};
				\node[above,font=\tiny] at (4,0.1) {$P_2$};
				\draw[<->,aerogreen,thick] (0,-0.3) -- (4,-0.3);
				\node[below,font=\scriptsize,aerogreen] at (2,-0.3)
				{$d_\mathrm{px}$};
				\node[below,font=\tiny] at (2,-0.7)
				{Barra de escala SEM};
			\end{tikzpicture}
		\end{column}
	\end{columns}
\end{frame}

\begin{frame}{Herramienta de Medición Interactiva}
	\begin{columns}
		\begin{column}{0.5\textwidth}
			\textbf{Notebook:} \texttt{CalculadoraDistancia}

			\vspace{0.3em}
			\begin{enumerate}\small
				\item Paso~A: Seleccionar imagen (dropdown)
				\item Paso~B: Calibrar escala
				      \begin{itemize}\scriptsize
					      \item Clic en extremos de la barra
					      \item Ingresar distancia real (\si{\micro\m})
					      \item Fijar escala
				      \end{itemize}
				\item Paso~C: Medir distancias
				      \begin{itemize}\scriptsize
					      \item Clic en dos puntos
					      \item Resultado en \si{\micro\m}
					      \item Repetir para 15+ mediciones
				      \end{itemize}
			\end{enumerate}

			\vspace{0.3em}
			\begin{block}{Salida}
				Cada medición reporta:\\
				$d_\mathrm{px}$ y $d_\mathrm{real}$ (\si{\micro\m})
			\end{block}
		\end{column}
		\begin{column}{0.47\textwidth}
			\begin{tikzpicture}[scale=0.7,>=Stealth,
					box/.style={rectangle,draw=aerogreen,thick,
							fill=aerogreen!8,text width=3.2cm,align=center,
							minimum height=0.6cm,font=\tiny,
							rounded corners=2pt}]
				\node[box] (sel) {Dropdown: seleccionar imagen};
				\node[box,below=0.5cm of sel]
				(cal) {Botón: Calibrar\\(2 clics + dist.)};
				\node[box,below=0.5cm of cal]
				(fix) {Botón: Fijar Escala};
				\node[box,below=0.5cm of fix]
				(meas) {Botón: Medir Distancia\\(2 clics)};
				\node[box,below=0.5cm of meas,draw=aerored,fill=aerored!8]
				(out) {Resultado: $d$ (\si{\micro\m})};
				\foreach \a/\b in {sel/cal,cal/fix,fix/meas,meas/out}
				\draw[->,aerogreen,thick] (\a) -- (\b);
				% Loop
				\draw[aerored,thick,->] (out.east) -- ++(0.8,0)
				|- (meas.east)
				node[pos=0.25,right,font=\tiny,aerored] {repetir};
			\end{tikzpicture}
		\end{column}
	\end{columns}
\end{frame}

% =============================================================================
\section{Evaluación de Riesgos}
% =============================================================================

\begin{frame}{Riesgos de la Práctica}
	\centering\small
	\begin{tabular}{>{\raggedright\arraybackslash}p{3.2cm}
		c c c c l}
		\toprule
		\textbf{Riesgo} & \textbf{P}                             & \textbf{S}
		                & $\boldsymbol{R}$                       & \textbf{Nivel}
		                & \textbf{Controles}                                          \\
		\midrule
		\rowcolor{aeroyellow!15}
		Alta tensión SEM
		                & 2                                      & 4              & 8
		                & \textcolor{aeroyellow}{\textbf{Medio}}
		                & Interlock, formación                                        \\
		\rowcolor{aeroyellow!15}
		Preparación metalo.
		                & 2                                      & 3              & 6
		                & \textcolor{aeroyellow}{\textbf{Medio}}
		                & Gafas, guantes, SOPs                                        \\
		\rowcolor{aeroyellow!15}
		Reactivos químicos
		                & 2                                      & 3              & 6
		                & \textcolor{aeroyellow}{\textbf{Medio}}
		                & Campana, MSDS                                               \\
		\rowcolor{aerogreen!15}
		Ergonomía digital
		                & 1                                      & 1              & 1
		                & \textcolor{aerogreen}{\textbf{Bajo}}
		                & Pausas, postura                                             \\
		\bottomrule
	\end{tabular}

	\vspace{0.4em}
	\begin{block}{$R = \text{Probabilidad} \times \text{Severidad}$}
		Bajo ($<5$) \quad | \quad Medio (5--12)
		\quad | \quad Alto ($>12$)
	\end{block}
\end{frame}

\begin{frame}{Matriz de Riesgos $5\times5$}
	\centering
	\begin{tikzpicture}[scale=0.65]
		\foreach \x in {0,...,4} {
				\foreach \y in {0,...,4} {
						\pgfmathsetmacro{\risk}{(\x+1)*(\y+1)}
						\pgfmathparse{\risk < 5 ? "aerogreen!40" :
							(\risk < 13 ? "aeroyellow!60" : "aerored!50")}
						\edef\fillcol{\pgfmathresult}
						\fill[\fillcol] (\x,\y) rectangle (\x+1,\y+1);
						\draw[gray] (\x,\y) rectangle (\x+1,\y+1);
						\node[font=\tiny] at (\x+0.5,\y+0.5)
						{\pgfmathprintnumber[fixed,precision=0]{\risk}};
					}
			}
		\foreach \x/\l in {0/MuyBajo,1/Bajo,2/Medio,3/Alto,4/MuyAlto} {
				\node[font=\tiny,rotate=45,anchor=west]
				at (\x+0.5,-0.1) {\l};
			}
		\foreach \y/\l in {0/Insignif.,1/Menor,2/Moder.,3/Mayor,
				4/Catastrófico} {
				\node[font=\tiny,anchor=east] at (-0.1,\y+0.5) {\l};
			}
		\node[font=\tiny\bfseries,rotate=90] at (-1.3,2.5)
		{Severidad};
		\node[font=\tiny\bfseries] at (2.5,-1) {Probabilidad};
		% Markers
		\fill[aerored,opacity=0.8] (1.5,3.5) circle (0.2);
		\node[white,font=\tiny\bfseries] at (1.5,3.5) {SEM};
		\fill[blue!70,opacity=0.8] (1.5,2.5) circle (0.2);
		\node[white,font=\tiny\bfseries] at (1.5,2.5) {Prep.};
		\fill[orange,opacity=0.8] (1.5,2.5) circle (0.2);
		\node[white,font=\tiny\bfseries] at (1.5,2.5) {Reac.};
		\fill[aerogreen!70,opacity=0.8] (0.5,0.5) circle (0.2);
		\node[white,font=\tiny\bfseries] at (0.5,0.5) {PC};
	\end{tikzpicture}

	\vspace{0.3cm}
	\begin{tikzpicture}[scale=0.7,every node/.style={font=\small}]
		\fill[aerogreendark] (0,0) -- (5,0) -- (4.2,0.8) -- (0.8,0.8) -- cycle;
		\fill[aerogreen] (0.8,0.8) -- (4.2,0.8) -- (3.4,1.6) -- (1.6,1.6) -- cycle;
		\fill[aeroyellow!80] (1.6,1.6) -- (3.4,1.6) -- (2.8,2.4) -- (2.2,2.4) -- cycle;
		\fill[aerored!60] (2.2,2.4) -- (2.8,2.4) -- (2.5,3.0) -- cycle;
		\node[white,font=\tiny\bfseries] at (2.5,0.4) {Eliminación / Sustitución};
		\node[white,font=\tiny\bfseries] at (2.5,1.2) {Ingeniería};
		\node[font=\tiny\bfseries] at (2.5,2.0) {Admin.};
		\node[white,font=\tiny\bfseries] at (2.5,2.6) {EPP};
		\draw[->,gray,very thick] (5.3,2.8) -- (5.3,-0.1);
		\node[gray,font=\tiny,rotate=90] at (5.6,1.3) {Más efectivo};
	\end{tikzpicture}
\end{frame}

% =============================================================================
\section{Casos de Estudio por Equipos}
% =============================================================================

\begin{frame}{Trabajo por Equipos: 4 Escenarios}
	\begin{block}{Instrucciones}
		Cada equipo recibe una imagen (o par de imágenes) diferente. Deben:
		\begin{enumerate}
			\item Ejecutar el pipeline completo sobre la(s) imagen(es) asignada(s)
			\item Calibrar la escala con la barra del SEM (notebook)
			\item Medir al menos 15 características dimensionales
			\item Diagnosticar el modo de fractura predominante
			\item Evaluar el riesgo: $P$ (1--5), $S$ (1--5),
			      $R = P \times S$, ubicar en la matriz
			\item Resolver 3 problemas adicionales
		\end{enumerate}
	\end{block}
\end{frame}

\begin{frame}{Equipo~1: Análisis Dimensional de Estrías}
	\begin{columns}
		\begin{column}{0.38\textwidth}
			\centering\scriptsize
			\begin{tabular}{ll}
				\toprule
				\textbf{Par.}          & \textbf{Valor}            \\
				\midrule
				Imagen                 & \texttt{fatigue3.png}     \\
				$N_\mathrm{reg}$       & 418                       \\
				Mediciones             & $\geq$15 estrías          \\
				Paris: $C$             & $6{,}9\!\times\!10^{-12}$ \\
				Paris: $m$             & 3                         \\
				$\Delta K_\mathrm{th}$ & $\SI{5}{\MPa\sqrt{\m}}$   \\
				\bottomrule
			\end{tabular}
		\end{column}
		\begin{column}{0.28\textwidth}
			\begin{alertblock}{\scriptsize Riesgo}
				\scriptsize
				Operación SEM (EHT),\\
				preparación metalográfica.\\
				Evaluar $P$, $S$, $R$.
			\end{alertblock}
		\end{column}
		\begin{column}{0.3\textwidth}
			\begin{exampleblock}{\scriptsize Problemas}
				\scriptsize
				\begin{enumerate}
					\item[P1.1] Error 5px en barra\\
					      200px $\to \varepsilon$?
					\item[P1.2] CLAHE clip 0,01\\
					      vs.\ 0,05
					\item[P1.3] Criterio $e$ para\\
					      estrías vs.\ hoyuelos
				\end{enumerate}
			\end{exampleblock}
		\end{column}
	\end{columns}
\end{frame}

\begin{frame}{Equipo~2: Comparación de Modos de Fractura}
	\begin{columns}
		\begin{column}{0.38\textwidth}
			\centering\scriptsize
			\begin{tabular}{ll}
				\toprule
				\textbf{Par.}    & \textbf{Valor}          \\
				\midrule
				Imágenes         & \texttt{fatigue4.png}   \\
				                 & \texttt{fatigue7.png}   \\
				$N_\mathrm{reg}$ & 153 / 549               \\
				Mediciones       & $\geq$15 por imagen     \\
				Diagnóstico      & dúctil / frágil / mixto \\
				\bottomrule
			\end{tabular}
		\end{column}
		\begin{column}{0.28\textwidth}
			\begin{alertblock}{\scriptsize Riesgo}
				\scriptsize
				Dos zonas del mismo\\
				componente.\\
				Comparar descriptores\\
				de forma.\\
				Evaluar $P$, $S$, $R$.
			\end{alertblock}
		\end{column}
		\begin{column}{0.3\textwidth}
			\begin{exampleblock}{\scriptsize Problemas}
				\scriptsize
				\begin{enumerate}
					\item[P2.1] Hoyuelos equiax.\\
					      vs.\ alargados
					\item[P2.2] Canny $\sigma$=1\\
					      vs.\ $\sigma$=3
					\item[P2.3] Detectar fractura\\
					      intergranular
				\end{enumerate}
			\end{exampleblock}
		\end{column}
	\end{columns}
\end{frame}

\begin{frame}{Equipo~3: Efecto de Resolución / Magnificación}
	\begin{columns}
		\begin{column}{0.38\textwidth}
			\centering\scriptsize
			\begin{tabular}{ll}
				\toprule
				\textbf{Par.}    & \textbf{Valor}                              \\
				\midrule
				Imágenes         & \texttt{fatigue5.png}                       \\
				                 & \texttt{fatigue4.png}                       \\
				$N_\mathrm{reg}$ & 2066 / 153                                  \\
				Tamaño           & \SI{2.3}{\mega\byte}\ / \SI{68}{\kilo\byte} \\
				Objetivo         & Cuantificar efecto                          \\
				\bottomrule
			\end{tabular}
		\end{column}
		\begin{column}{0.28\textwidth}
			\begin{alertblock}{\scriptsize Riesgo}
				\scriptsize
				Alta vs.\ baja resolución.\\
				¿Son comparables las\\
				mediciones absolutas?\\
				Evaluar $P$, $S$, $R$.
			\end{alertblock}
		\end{column}
		\begin{column}{0.3\textwidth}
			\begin{exampleblock}{\scriptsize Problemas}
				\scriptsize
				\begin{enumerate}
					\item[P3.1] Menor caract.\\
					      medible (3px mín.)
					\item[P3.2] Resize 25\% y\\
					      reprocesar
					\item[P3.3] Protocolo para\\
					      estadísticas repr.
				\end{enumerate}
			\end{exampleblock}
		\end{column}
	\end{columns}
\end{frame}

\begin{frame}{Equipo~4: Optimización de Parámetros}
	\begin{columns}
		\begin{column}{0.38\textwidth}
			\centering\scriptsize
			\begin{tabular}{ll}
				\toprule
				\textbf{Par.}           & \textbf{Valor}                       \\
				\midrule
				Imagen                  & \texttt{fatigue6.png}                \\
				$N_\mathrm{reg}$ (base) & 2897                                 \\
				clip\_limit             & $\{0{,}01; 0{,}03; 0{,}05; 0{,}10\}$ \\
				$\sigma_\mathrm{Canny}$ & $\{1; 2; 3; 5\}$                     \\
				disk                    & $\{1; 2; 3; 5\}$                     \\
				\bottomrule
			\end{tabular}
		\end{column}
		\begin{column}{0.28\textwidth}
			\begin{alertblock}{\scriptsize Riesgo}
				\scriptsize
				Imagen más compleja.\\
				Optimizar pipeline\\
				para maximizar calidad\\
				de segmentación.\\
				Evaluar $P$, $S$, $R$.
			\end{alertblock}
		\end{column}
		\begin{column}{0.3\textwidth}
			\begin{exampleblock}{\scriptsize Problemas}
				\scriptsize
				\begin{enumerate}
					\item[P4.1] disk que elimine\\
					      todas las regiones
					\item[P4.2] Opening + Closing\\
					      ¿efecto?
					\item[P4.3] Otsu vs.\ umbral\\
					      fijo $t=128$
				\end{enumerate}
			\end{exampleblock}
		\end{column}
	\end{columns}
\end{frame}

% =============================================================================
\section{Herramientas de Software}
% =============================================================================

\begin{frame}[fragile]{Script \texttt{P05\_Analisis\_Imagen.py}}
	\begin{columns}
		\begin{column}{0.5\textwidth}
			\begin{block}{Ruta}
				\scriptsize
				\texttt{Practicas/P05\_Analisis\_Imagen/}\\
				\texttt{\ \ src/P05\_Analisis\_Imagen.py}
			\end{block}
			\vspace{0.2em}
			\scriptsize
			\begin{tabular}{ll}
				\toprule
				\textbf{Función}
				 & \textbf{Resultado}          \\
				\midrule
				\texttt{cargar\_imagen\_gris}
				 & float $[0,1]$               \\
				\texttt{procesar\_clahe}
				 & Imagen CLAHE                \\
				\texttt{procesar\_sobel}
				 & Mapa gradientes             \\
				\texttt{procesar\_canny}
				 & Bordes binarios             \\
				\texttt{procesar\_otsu}
				 & Máscara binaria             \\
				\texttt{procesar\_morfologia}
				 & Máscara limpia              \\
				\texttt{procesar\_etiquetado}
				 & $N$ regiones + descriptores \\
				\texttt{run\_full\_analysis}
				 & Pipeline + 6 PNG            \\
				\bottomrule
			\end{tabular}
		\end{column}
		\begin{column}{0.47\textwidth}
			\begin{lstlisting}[language=Python]
# Ejecucion desde la raiz:
python Practicas\P05_Analisis_Imagen\
       src\P05_Analisis_Imagen.py

# Salida: 30 archivos PNG
# (6 por imagen x 5 imagenes)
# en data/images/

# Patron de nombres:
#   fatigue3_00_resumen.png
#   fatigue3_01_clahe.png
#   fatigue3_02_sobel.png
#   fatigue3_03_canny.png
#   fatigue3_04_otsu.png
#   fatigue3_05_opening.png
			\end{lstlisting}

			\vspace{0.2em}
			\begin{exampleblock}{Notebook}
				\scriptsize
				Misma lógica + calibración interactiva\\
				+ medición con clics en la imagen
			\end{exampleblock}
		\end{column}
	\end{columns}
\end{frame}

\begin{frame}[fragile]{Ejemplo de Uso del Pipeline}
	\begin{lstlisting}[language=Python]
# Importar funciones del pipeline
from P05_Analisis_Imagen import *

# 1. Cargar imagen en escala de gris
img = cargar_imagen_gris('files/p5/fatigue3.png')

# 2. Mejorar contraste
img_clahe = procesar_clahe(img, clip_limit=0.03)
img_u8 = img_as_ubyte(img_clahe)

# 3. Detectar bordes
img_sobel = procesar_sobel(img_u8)
img_canny = procesar_canny(img_u8, sigma=2.0)

# 4. Segmentar
img_otsu = procesar_otsu(img_u8)

# 5. Limpiar morfologicamente
img_clean = procesar_morfologia(img_otsu, 'opening', 2)

# 6. Etiquetar y contar regiones
label_img, props = procesar_etiquetado(img_clean)
# >> Numero de objetos encontrados: 418
	\end{lstlisting}
\end{frame}

% =============================================================================
\section{Cuestionario y Formato}
% =============================================================================

\begin{frame}{Preguntas de Discusión}
	\begin{enumerate}\small
		\item ¿Por qué la calibración de escala es el paso más crítico
		      para mediciones cuantitativas?
		\item Diferencia entre Sobel y Canny.\ ¿Cuándo usar cada uno?
		\item Escriba $\sigma_B^2(t)$ de Otsu y explique por qué
		      maximizarla es equivalente a minimizar la varianza
		      intra-clase.
		\item ¿Por qué el SEM requiere muestra conductora?
		      ¿Qué artefactos aparecen si no lo es?
		\item Demuestre que $I \circ SE \subseteq I$
		      (opening nunca agrega píxeles).
		\item ¿Qué infiere de hoyuelos pequeños y equiaxiales vs.\
		      grandes y alargados?
		\item ¿Por qué es importante analizar a múltiples
		      magnificaciones?
		\item Efecto del EHT en el contraste: ¿alto o bajo EHT para
		      fractografía superficial?
		\item Si no hay estrías claras pero se sospecha fatiga,
		      ¿qué buscar?
		\item Proponga un método cuantitativo para clasificar
		      ``dúctil'' vs.\ ``frágil'' con descriptores de región.
	\end{enumerate}
\end{frame}

\begin{frame}{Formato del Informe (AIAA)}
	\begin{enumerate}
		\item \textbf{Resumen:} objetivo, pipeline, resultados
		\item \textbf{Introducción:} importancia fractografía SEM
		\item \textbf{Fundamento Teórico:} SEM, fractura, algoritmos
		\item \textbf{Metodología:} pipeline, parámetros, calibración
		\item \textbf{Resultados:} 6 etapas del pipeline + tabla de
		      regiones + mediciones + evaluación de riesgo
		\item \textbf{Discusión:} diagnóstico + cuestionario + efecto
		      de parámetros
		\item \textbf{Conclusiones:} 4--6 puntos
		\item \textbf{Referencias} (IEEE)
		\item \textbf{Apéndice~A:} imágenes del script/notebook
	\end{enumerate}
\end{frame}

% =============================================================================
\section{Conclusiones}
% =============================================================================

\begin{frame}{Conclusiones}
	\begin{enumerate}\small
		\item El pipeline de 6 etapas permite extraer información
		      \textbf{cuantitativa} de micrografías SEM
		\item CLAHE es esencial para iluminación desigual
		      (controlado por \texttt{clip\_limit})
		\item Otsu: umbral \textbf{óptimo y automático},
		      sin subjetividad
		\item Operaciones morfológicas: indispensables para
		      eliminar ruido residual
		\item Regiones detectadas (153--2897) dependen de la
		      resolución $\to$ reportar condiciones de adquisición
		\item Calibración de escala: paso más crítico
		      (5~px error $\to$ $\sim\!5\%$ error en mediciones)
		\item Script automatizado: 5 imágenes $\to$ 30 archivos,
		      reproducible y trazable
	\end{enumerate}
\end{frame}

\begin{frame}{Resumen Visual: Pipeline Completo}
	\centering
	\includegraphics[width=\textwidth,height=0.62\textheight,keepaspectratio]{fatigue3_00_resumen.png}

	\vspace{0.2cm}
	\begin{tikzpicture}
		\node[draw=aerogreen,very thick,rounded corners,fill=aerogreen!10,
			text width=10cm,align=center,font=\small\bfseries] {
			De la micrografía SEM al diagnóstico cuantitativo:\\
			6 etapas automatizadas + calibración + medición
		};
	\end{tikzpicture}
\end{frame}

\begin{frame}{}
	\centering
	\vspace{2cm}
	{\Huge\color{aerogreen}\bfseries ¿Preguntas?}\\[1cm]
	{\large P05 -- Análisis de Imagen Digital y Fractografía SEM}\\[0.5cm]
	{\small Técnicas de Medida -- Ingeniería Aeronáutica}
\end{frame}

\end{document}
